\thesistranslationchinese

\section{摘要}

本文针对将英语句子映射为Bash命令（NL2Bash）的问题，提出了新的数据和语义解析方法。我们的长期目标是让用户仅通过英语描述即可完成文件操作、搜索和应用脚本等任务。作为该领域的第一步，我们提供了一个包含常用但复杂的Bash命令及专家撰写的英语描述的新数据集，并建立了基准方法以确定该任务的性能水平。

\textbf{关键词：} 自然语言编程、自然语言接口、语义解析

\section{引言}

用英语或其他自然语言编程的梦想几乎与编程任务本身一样古老（Sammet, 1966）。尽管自然语言的精确性远不如形式语言（Dijkstra, 1978），但作为编程媒介，自然语言具有普适性，并支持高度重复任务的自动化，如文件操作、搜索和应用脚本（Wilensky等, 1984; Wilensky等, 1988; Dahl等, 1994; Quirk等, 2015; Desai等, 2016）。

本研究在一个新颖且富有挑战性的领域——操作系统的自然语言控制——提出了新的数据和语义解析方法。我们的长期目标是让用户仅通过自然语言（NL）描述即可在计算机上执行任务。作为该方向的第一步，我们提供了一个大规模的新数据集（NL2Bash），包含常用但复杂的命令和专家撰写的描述，并建立了基准方法以确定该任务的性能水平。

NL2Bash问题可视为一种语义解析，其目标是将句子映射到其底层含义的形式化表示（Mooney, 2014）。我们引入的数据集提供了一种新的目标意义表示类型（Bash命令），并且比大多数现有的语义解析基准数据集大两到十倍（Dahl等, 1994; Popescu等, 2003）。

其他近期的语义解析工作也聚焦于编程语言，包括正则表达式（Locascio等, 2016）、IFTTT脚本（Quirk等, 2015）和SQL查询（Kwiatkowski等, 2013; Iyer等, 2017; Zhong等, 2017）。然而，我们所考虑的Shell命令数据带来了独特的挑战，原因是其不规则的语法（命令选项没有语法树表示）、广泛的领域覆盖（>100个Bash工具）以及大量的未登录词（例如命令可以操作任意文件）。

我们通过从问答论坛、教程、技术博客和课程材料等网站抓取常用Bash命令来构建NL2Bash语料库。我们从Bash程序员那里收集了一组高质量的命令描述。经过严格的质量控制，我们收集了超过9,000个英语-命令对，涵盖了100多个独特的Bash工具。

我们还通过一系列实验证明NL2Bash是一个值得未来研究的挑战性任务。我们基于近期的神经语义解析工作（Dong和Lapata, 2016; Ling等, 2016），评估了标准的Seq2seq模型（Sutskever等, 2014）和CopyNet模型（Gu等, 2016）。我们还纳入了一个最近提出的分阶段神经语义解析模型Tellina，该模型使用手动定义的启发式规则来更好地检测和翻译命令参数（Lin等, 2017）。

我们发现，当在正确的序列粒度（子词元）上应用时，CopyNet显著优于分阶段模型，且需要的预处理和后处理明显更少。我们表现最佳的系统在top-1命令结构准确率达到49\%，top-1完整命令准确率达到36\%。这些性能水平虽然远非完美，但在精心设计的界面中已足够实用（Lin等, 2017），同时也表明未来建模创新仍有充足空间。

\section{领域：Linux Shell命令}

Shell命令由三个基本组件组成：工具（如find、grep）、选项标志（如-name、-i）和参数（如"*.java"、"TODO"）。工具可以有惯用的标志语法。

Bash工具有250多个，第三方开发者还会定期添加新工具。我们专注于Linux用户小组确定的135个最有用的工具。我们只考虑可以在单行（one-liners）中指定的目标命令。

我们省略了包含I/O重定向、变量赋值和复合语句等语法结构的命令，因为这些命令需要在上下文中进行解释。

\section{语料库构建}

语料库由文本-命令对组成，每个对包含一个从网络抓取的Bash命令和一个专家生成的自然语言描述。

\subsection{数据收集}

我们雇用了10名熟悉Shell脚本的Upwork自由职业者。他们从问答论坛、教程、技术博客和课程材料等网页中收集文本-命令对。

自由职业者从网页中复制Bash命令，并根据其背景知识和网页上下文复制或编写文本。我们将自然语言描述限制为单个句子，Bash命令限制为单行命令。

自由职业者为网页上的每个命令提供一个自然语言描述。一个自由职业者可能在不同网页中多次标注同一命令，多个自由职业者也可能标注同一命令。收集多个不同的描述增加了数据集的语言多样性。

\subsection{数据清洗}

我们使用自动化流程来过滤和清洗数据集。

清洗脚本移除了：

\begin{itemize}
    \item 违反Linux手册页语法规范的非语法命令。
    \item 包含超出范围语法结构的命令。
    \item 主要用于多语句Shell脚本的命令。
    \item 包含非Bash语言解释器的命令。
\end{itemize}

我们使用概率拼写检查器纠正了自然语言描述中的拼写错误。我们还手动纠正了拼写检查器未能捕获的自然语言和Shell命令中的部分拼写错误。

对于Bash命令，我们从每个命令开头移除了sudo以及Shell输入提示字符（如"\$"和"\#"）。我们将工具的绝对路径名替换为其基本名称。

\subsection{语料库统计}

经过过滤和清洗，我们的数据集包含9,305个对。Bash命令涵盖102个独特的工具，使用了206个标志——这是一个功能丰富的领域。

自然语言句子和Bash命令的平均长度相对较短，分别为11.7个词和7.7个词元。词频和命令词元频率的中位数均为1，这是由大量开放词汇常量造成的。

我们将命令模板定义为将其参数替换为语义类型的命令。

虽然大多数自然语言句子和Bash命令形成一对一映射，但该问题本质上是一个多对多映射问题——存在许多语义等价的命令，且一个Bash命令可能用不同的自然语言描述来表达。

\subsection{数据划分}

我们将过滤后的数据划分为训练集、开发集和测试集。

我们首先按自然语言描述对对进行聚类。我们以10:1:1的比例将聚类随机划分为训练集、开发集和测试集。

划分后，我们将所有命令出现在训练集中的开发集和测试集对移至训练集。这防止了模型通过简单地记忆训练集中见过的自然语言描述或命令来获得高准确率。

\section{评估方法}

在我们的数据集中，一个自然语言描述可能有多个正确的Bash命令翻译。这给评估带来了挑战，因为并非所有正确的命令都出现在我们的数据集中。

我们雇用了3名熟悉Shell脚本的Upwork自由职业者。为了评估特定系统，自由职业者独立评估了其对所有测试样本的top-3翻译的正确性。

我们报告两种类型的准确率：

\begin{itemize}
    \item Top-k完整命令准确率
    \item Top-k命令模板准确率
\end{itemize}

我们将完整命令准确率定义为模型输出中正确完整命令排名在k或以上的测试实例百分比。

我们将命令模板准确率定义为模型输出中正确命令模板排名在k或以上的测试实例百分比。

NL到代码翻译的先前方法也注意到了类似的问题。一些方法形式化验证等价性，而另一些方法则在虚拟环境中执行生成的代码并比较结果。

论文认为BLEU在衡量形式语言的语义相似度方面效果不佳。

\section{系统设计挑战}

本节列出了Bash领域语义解析的挑战。

\subsection{丰富的领域}

Bash的应用领域从文件系统管理、文本处理、网络控制到高级操作系统功能（如进程管理）。

\subsection{未登录常量}

Bash命令包含许多开放词汇常量，如文件/路径名、文件属性和时间表达式。

\subsection{语言灵活性}

许多Bash命令有大量选项标志，多个命令可以组合以解决更复杂的任务。

\subsection{惯用语法}

Bash解释器使用浅层语法文法来解析管道、代码块和其他高级语法结构。每个命令可以有自己的惯用语法规则。

\section{基准系统性能}

为了为未来工作建立性能水平，我们评估了两种神经机器翻译模型：Seq2Seq和CopyNet。我们还评估了Tellina。

\subsection{Seq2Seq}

Seq2Seq模型使用RNN编码器-解码器定义给定输入序列的输出序列的条件概率。

\subsection{CopyNet}

CopyNet是Seq2Seq的扩展，能够在生成输出序列时选择输入序列的子序列并在适当位置输出。

\subsection{Tellina}

Tellina首先将自然语言中的常量抽象为相应的语义类型，并执行模板级别的NL到代码翻译。然后填充代码模板中的参数槽位。

\subsection{实现细节}

我们使用门控循环单元（GRU）RNN单元和双向RNN编码器。

我们在三个词元粒度级别上评估了Seq2Seq和CopyNet：

\begin{itemize}
    \item 词元级别
    \item 字符级别
    \item 子词元级别
\end{itemize}

为了计算子词元，我们将自然语言和Bash命令中的每个常量拆分为字母和数字的连续序列。

我们使用小批量Adam优化学习目标。

\subsection{结果}

实验表明：

\begin{itemize}
    \item 词元级别建模产生更高的命令结构准确率。
    \item 字符和子词元模型给出更高的完整命令准确率。
    \item CopyNet显著提高了完整命令准确率。
    \item 子词元CopyNet达到了最高的完整命令准确率。
\end{itemize}

Tellina模型产生了第二好的完整命令准确率和第二好的结构准确率。

\subsection{错误分析}

作者手动检查了top-1系统输出。

主要错误原因包括：

\begin{itemize}
    \item 训练数据稀疏性
    \item 工具错误
    \item 标志错误
    \item 参数格式错误
    \item 常量枚举
    \item 复杂任务
\end{itemize}

论文建议未来工作在解码器中构建浅层命令结构，而不是以纯序列方式合成整个输出。

\section{与现有数据集的比较}

论文将NL2Bash与其他语义解析和NL到代码数据集进行了比较。

大多数现有数据集是为领域特定语言构建的。一些近期数据集使用Java、Python、C\#和Bash。

NL2Bash是使用实际代码片段和专家撰写的自然语言构建的最大数据集之一。

论文讨论了多种数据收集策略：

\begin{itemize}
    \item 人工标注
    \item 自动抓取
    \item 合成生成
\end{itemize}

作者认为这些策略是互补的。

\section{结论}

我们通过引入一个大规模的新数据集和基准方法，研究了将英语句子映射为Bash命令（NL2Bash）的问题。

NL2Bash是使用实际代码片段和专家撰写的自然语言构建的最大NL到代码数据集之一。

实验证明了现有模型的竞争性能，同时也表明在这个具有挑战性的语义解析问题上，未来工作仍有充足空间。
